% SPDMX — ICASSP proceedings outline
% This file contains narrative notes only. Replace each note block with prose.
% --------------------------------------------------------------------------
\documentclass{article}
\usepackage{spconf,hyperref}

\newcommand{\demourl}{https://pnlong.github.io/SPDMX/}

\title{SPDMX: Public-Domain Multitrack Stem Audio at Scale}

\name{Phillip Long, Eduardo Escoto, Julian McAuley, Zachary Novack}
\address{University of California, San Diego Computer Science and Engineering Department}

\begin{document}
%\ninept
\maketitle

% ABSTRACT — one sentence or compact clause for each move:
% - Problem: legally redistributable multitrack audio is scarce; existing
%   resources trade off realism, scale, and redistributability.
% - Opportunity and obstacle: public-domain symbolic music offers scale, but
%   naive rendering produces acoustically weak stems.
% - Solution: introduce SPDMX, a listening-guided, fully redistributable
%   rendering of PDMX with aligned scores, MIDI, stems, and mixtures.
% - Method: correct score metadata, select category-specific open renderers
%   through a listening study, and preserve musical balance during mixing.
% - Scale: 253,982 songs, 5,469 song-hours, and 490,159 mono stems; distinguish
%   clearly between the full corpus and the 88,403 songs with two or more stems.
% - Evidence: matched-budget source-separation experiments test in-domain
%   learnability, complementarity with Slakh2100, and transfer to recorded
%   datasets including MUSDB18, MoisesDB, and MedleyDB.
% - Takeaway: SPDMX makes large-scale stem supervision reproducible and
%   redistributable from source score to final mixture.

% KEYWORD CANDIDATES:
% Public-domain music; audio datasets; multitrack stems; symbolic-to-audio;
% music source separation.

\section{Introduction}
\label{sec:intro}

% PARAGRAPH 1 — NEED
% Topic sentence: Modern music models require paired mixtures and stems, but
% legally redistributable multitrack audio remains scarce.
% - Name motivating tasks: separation, accompaniment generation, automatic
%   mixing, and stem-conditioned synthesis.
% - Briefly contrast music with fields that have massive public corpora.

% PARAGRAPH 2 — THE EXISTING TRADE-OFF
% Topic sentence: Existing resources trade off realism, scale, and
% redistributability.
% - Recorded datasets (MUSDB18, MedleyDB, MoisesDB): realistic but small and/or
%   restricted.
% - Slakh2100: synthetic and larger, but built with non-redistributable
%   commercial instruments and still small relative to symbolic collections.
% - State the gap explicitly: no existing resource combines all three goals.
% - Place the dataset comparison table here.

% PARAGRAPH 3 — OPPORTUNITY AND TECHNICAL OBSTACLE
% Topic sentence: Public-domain symbolic collections offer a path to scale, but
% naive rendering produces acoustically unconvincing stems.
% - Briefly explain PDMX's crowdsourced provenance, public-domain status,
%   contents, and scale for readers unfamiliar with it.
% - Explain why direct General MIDI rendering is insufficient: dry timbres and
%   poor sustain for winds, brass, and strings.

% PARAGRAPH 4 — SPDMX
% Topic sentence: We introduce SPDMX, a listening-guided rendering of PDMX that
% pairs public-domain scores with redistributable stems and mixtures at
% unprecedented scale.
% - Give headline corpus statistics.
% - Distinguish all rendered songs from the multitrack subset.
% - Emphasize the complete open chain: source score, MIDI, synthesizers,
%   rendering protocol, stems, mixtures, and metadata.
% - Use "to our knowledge" once and support the largest-dataset claim with the
%   comparison table.

% PARAGRAPH 5 — CONTRIBUTIONS AND PAPER PROMISE
% Topic sentence: SPDMX combines a fully open synthesis protocol,
% listener-selected rendering strategies, and task-based evaluations.
% - Contribution 1: aligned symbolic and audio corpus at public-domain scale.
% - Contribution 2: category-specific hybrid FluidSynth/MIDI-DDSP recipe.
% - Contribution 3: perceptual study used to freeze that recipe.
% - Contribution 4: matched-budget source-separation evaluation spanning
%   synthetic and recorded test domains.
% - End with release status and project-page link; avoid making the project
%   page carry evidence needed to understand the paper.

\section{Related work}
\label{sec:prior}

% PARAGRAPH 1 — RECORDED MULTITRACK DATA
% Topic sentence: Recorded multitrack datasets provide realistic supervision
% but remain too small and restricted for large-scale open training.
% - Situate MUSDB18, MedleyDB, and MoisesDB.
% - Explain what realism provides and what scale/licensing prevents.

% PARAGRAPH 2 — SYNTHETIC STEM DATA
% Topic sentence: Synthetic datasets increase scale, but prior large
% collections depend on proprietary sample libraries.
% - Present Slakh2100 as the closest precedent.
% - Explain what SPDMX retains: coherent stems and mixtures from symbolic music.
% - Explain what SPDMX changes: public-domain scores, open renderers, much
%   greater scale, and end-to-end redistributability.

% PARAGRAPH 3 — NEURAL MIDI SYNTHESIS
% Topic sentence: Neural MIDI synthesis can improve realism for selected
% instruments, motivating a hybrid rather than single-renderer pipeline.
% - Discuss MIDI-DDSP, URMP, and the Chamber Ensemble Generator.
% - Set up the listening study: renderer choice is empirical and
%   instrument-dependent, not an assumption that neural synthesis always wins.

\section{SPDMX Dataset}
\label{sec:dataset}

% SECTION FRAMING
% Topic sentence: Constructing an open stem corpus from symbolic scores
% requires solving three problems: unreliable instrumentation metadata,
% uneven rendering quality, and mixture assembly that preserves musical
% balance.
% - Map these problems directly to the three subsections below.

\subsection{Initial data cleaning}
\label{ssec:cleaning}

% PARAGRAPH 1 — WHY CLEANING IS NECESSARY
% Topic sentence: We first correct instrument metadata and remove empty tracks
% so that each rendered stem has a reliable semantic identity.
% - Explain disagreement between track names and General MIDI programs.
% - Describe alias-based program correction and removal of tracks without notes.
% - State limitations of heuristic correction if relevant.

% PARAGRAPH 2 — WHAT THE CLEANED INVENTORY CONTAINS
% Topic sentence: The cleaned corpus is dominated by piano, voice/choir,
% drums, and common orchestral families, with different distributions by stem
% count and audible duration.
% - Define non-silent hours and the RMS threshold.
% - Explain why piano's share of stems exceeds its share of non-silent hours.
% - Place the corrected-program inventory figure here.

% TRANSITION
% Topic sentence: Cleaning makes the symbolic controls reliable, but it does
% not resolve the acoustic limitations of naive single-bank rendering.
% - Motivate the category-specific listening comparison.

\subsection{Listening study}
\label{ssec:listening}

% PARAGRAPH 1 — DESIGN SPACE
% Topic sentence: Because no open synthesizer performs best across all
% instrument families, we compare category-specific FluidSynth and MIDI-DDSP
% rendering strategies.
% - Define basic and varied FluidSynth conditions.
% - List representative redistributable soundfonts and the light reverb setup.
% - Define MIDI-DDSP eligibility and FluidSynth fallbacks.
% - Clarify treatment of voice tracks and exclusion of lyric synthesis.

% PARAGRAPH 2 — STUDY DESIGN
% Topic sentence: Listeners rate four rendering families for perceived realism
% under shuffled, instrument-balanced comparisons.
% - Four families: basic; varied; MIDI-DDSP with basic fallback; MIDI-DDSP with
%   varied fallback.
% - State excerpt length, number of categories and stems, listener count,
%   rating scale, reference handling, and randomization.
% - Be precise about the unit of analysis and how scores are aggregated.

% PARAGRAPH 3 — DECISION RULE AND RESULT
% Topic sentence: Listener preferences support a hybrid recipe whose renderer
% depends on instrument family.
% - State the default highest-mean rule and the diversity exception.
% - Give the frozen assignment by category.
% - Place the listening-study figure here.
% - Avoid implying statistical significance unless it is tested and reported.

\subsection{Final assembly}
\label{ssec:assembly}

% PARAGRAPH 1 — MIXING PROBLEM AND SOLUTION
% Topic sentence: We preserve MIDI velocity relationships while controlling
% loudness and clipping so that released stems remain linearly summable.
% - Explain why independent loudness normalization would erase intended balance.
% - Describe the -23 LUFS target, peak limiting, velocity-ratio scaling, and
%   one uniform anti-clipping gain on the stem sum.
% - State explicitly that each mixture equals the sum of its released stems.

% PARAGRAPH 2 — RELEASE CONTENTS AND SCALE
% Topic sentence: The resulting corpus provides aligned scores, MIDI, stems,
% and mixtures for 253,982 songs, including 88,403 multitrack songs.
% - Report 5,469 total song-hours, 490,159 mono FLAC stems at 44.1 kHz, and
%   2,734 multitrack song-hours.
% - Define song-hours and explain why they are not multiplied by stem count.
% - Report stem-count distribution and identify the BDGP comparison slice.
% - Explain identity-based alignment with PDMX metadata and MusicXML.
% - Name the open synthesis components and their redistribution terms.
% - Place the stems-per-song figure here.

\section{Experiments}
\label{sec:experiments}

% SECTION FRAMING
% Topic sentence: We evaluate whether SPDMX provides learnable supervision for
% source separation and whether it generalizes across synthetic and recorded
% datasets under matched training budgets.
% - State that this is a dataset validation study, not a claim of
%   state-of-the-art separation performance.
% - Explain exactly what "matched" controls: steps, examples, duration, or
%   another budget.

\subsection{Source separation}
\label{ssec:separation}

% PARAGRAPH 1 — RESEARCH QUESTION AND SETUP
% Topic sentence: We first test whether SPDMX supports in-domain source
% separation, complements Slakh across synthetic domains, and transfers to
% recorded multitrack datasets.
% - Define Hybrid Demucs and the bass/drums/guitar/piano target classes.
% - Explain SPDMX-to-class mapping and the all-four-stems eligibility filter.
% - Report training-set hours and the three matched-budget arms.
% - Define SI-SDR, aggregation, and uncertainty/repeated runs.
% - Identify the evaluation domains: held-out SPDMX, Slakh2100, MUSDB18,
%   MoisesDB, and MedleyDB.
% - Explain how each dataset's stem taxonomy is mapped to common target classes;
%   report results only for classes that can be compared fairly.

% PARAGRAPH 2 — PRIMARY RESULT
% Topic sentence: SPDMX provides strong in-domain supervision, while joint
% training retains much of Slakh performance and substantially improves
% transfer to SPDMX.
% - Report the in-domain and union results.
% - Interpret poor cross-domain performance as evidence of acoustic mismatch.
% - Keep claims proportional to the observed differences.
% - Place the separation table here.

% PARAGRAPH 3 — CROSS-DATASET GENERALIZATION
% Topic sentence: We next evaluate whether models trained on synthetic stems
% transfer to recorded music across MUSDB18, MoisesDB, and MedleyDB.
% - Present each external dataset separately or with per-dataset columns; do
%   not average across datasets with different taxonomies or distributions.
% - For the completed MUSDB18 comparison, report the weak transfer of all
%   synthetic training arms and note SPDMX's best point estimate cautiously.
% - Add MoisesDB and MedleyDB results once available; until then, describe them
%   as planned evaluations rather than evidence.
% - Use repeated runs or uncertainty before interpreting small differences.
% - If supported, discuss timbral diversity, rendering artifacts, class
%   balance, and dataset taxonomy as possible causes of transfer differences.
% - State the broader implication: openness and scale do not by themselves
%   eliminate the synthetic-to-recorded domain gap.
% - Place the cross-dataset separation table here.

\section{Conclusion}
\label{sec:conclusion}

% PARAGRAPH 1 — RETURN TO THE CENTRAL TRADE-OFF
% Topic sentence: SPDMX addresses the scale--redistributability gap in
% multitrack music data without relying on copyrighted recordings or
% proprietary instruments.
% - Briefly restate the listening-guided hybrid pipeline and release scale.
% - Avoid repeating every implementation detail or metric.

% PARAGRAPH 2 — WHAT THE EVIDENCE ESTABLISHES
% Topic sentence: Matched-budget source-separation experiments test whether
% SPDMX is learnable in-domain, complements Slakh, and transfers across recorded
% evaluation datasets.
% - Keep the conclusion conservative: validation and utility, not superiority.
% - Summarize MoisesDB and MedleyDB only after those results are available.
% - Treat synthetic-to-recorded transfer as a limitation if the broader
%   evaluation confirms the current MUSDB18 result.

% PARAGRAPH 3 — CAPABILITY ENABLED
% Topic sentence: By jointly releasing symbolic scores, MIDI, stems, and
% mixtures under redistributable terms, SPDMX enables music-model research at a
% scale previously available only through restricted data and proprietary
% synthesis.
% - Name likely uses: separation, stem-conditioned generation, automatic
%   mixing, transcription, accompaniment, and joint symbolic/audio modeling.
% - Give the project-page URL and concrete release status.
% - Mention transcription and accompaniment pilots only as future work unless
%   completed results are included; omit pending-result tables.

\section{Compliance with Ethical Standards}
\label{sec:ethics}

% PARAGRAPH 1 — PROVENANCE AND REDISTRIBUTION
% Topic sentence: SPDMX is derived from public-domain PDMX and uses only
% redistributable synthesis components.
% - Document score provenance, soundfont licenses, MIDI-DDSP weights, and any
%   release restrictions or attribution requirements.
% - Contrast the pipeline with closed commercial VST sample libraries.
% - Note relevant limitations or potential misuse if required by the venue.

\end{document}
